% =========================================================================
\chapter{Background and Related Work}
\label{ch:related}
% =========================================================================



% -------------------------------------------------------------------------
\section{Self-evolving agent harnesses}
\label{sec:self-evolving}
% -------------------------------------------------------------------------

The coordinates for this line are set by a recent survey~\cite{survey-selfevolve},
which organizes the field along what, when, how and where an agent evolves.
On its \emph{where} axis the object of this thesis is harness
configuration: not model weights, and not memory alone. This review covers
five relevant literature streams, including 2026 preprints; it is not
systematic, and claims of absence are limited to the covered literature,
with reinforcement-learning methods included only where they bear directly
on these mechanisms.

\paragraph{What the candidate is.} Two answers dominate. Free code:
ADAS~\cite{adas} lets an agent rewrite itself, and the Darwin G\"odel
Machine~\cite{dgm} adds a population archive of self-modifying agents ---
the nearest precedent for the layered hash identity of
\S\ref{sec:identity}. Or a workflow: AFlow~\cite{aflow} evolves code-represented agentic workflows with named nodes and operators, and GEPA~\cite{gepa} performs reflective prompt evolution.
AlphaEvolve~\cite{alphaevolve} contributes a cascade of increasingly demanding evaluations that filters faulty or unpromising programs early, using user-provided tests, scores and optional criteria. An in-run redundancy guard is another methodological
precedent~\cite{ledger-t17}.

\paragraph{A directly related problem setting.} HarnessForge~\cite{harnessforge}
targets the same object as this thesis: it treats the harness as an executable code/configuration bundle over
Planning, Action and Memory and evolves it by fault-attribution-guided archive
search with a trained policy. Candidates pass interface and schema checks
before Pareto filtering. It addresses the same object through a different
design, but has no graph-hash identity, native execution graph, or
reconciliation between what a candidate declares and what execution observed --- and it trains, where GHX is
training-free and structural.

\paragraph{The harness-level line in 2026.} The object of this thesis
acquired a literature of its own during the year it was written, and the
positions worth naming are these. \textbf{NLAH}~\cite{nlah} lifts the
harness out of controller code into an editable natural-language document
interpreted at runtime, cutting harness tokens by $95\%$ and beating a
code harness on Live-SWE; its own appendix concedes that natural-language
imprecision is ``under-specified, interpreted differently across models,
or weakened by paraphrase'' --- there is no machine-checkable constraint.
It is the representational rival of this thesis --- same problem, opposite
remedy --- and the conflict is taken up in \S\ref{sec:positioning}.
\textbf{SIGIL}~\cite{sigil} goes the other way, compiling prose skills
into a typed harness behind a compile gate that refuses any rule without
a source citation --- the nearest precedent for a build-time gate, on a
single skill rather than a whole harness. \textbf{HarnessBank}~\cite{harnessbank}
keeps a gene bank of harnesses and screens mutations before expensive
evaluation with a paired significance gate --- $\hat{\Delta} > 0$ and
$z \ge 1.96$, the critical value of a two-sided $5\%$ test applied in the
positive direction only --- reporting that it eliminates the \emph{false
elites} an ungated bank admits, gains indistinguishable from noise.
It is the one system in this survey that combines an explicit activation
check with a paired statistical gate rather than admitting on raw score, and it names non-regression criteria among the
screening rules it argues are too weak. Its seven domains include
SWE-bench, where its own report concedes that $n{=}26$ cannot separate the
arms; truncation is not treated.
\textbf{AgentDevel}~\cite{agentdevel} reframes self-evolution as release
engineering: versioned specifications, symptom-level diagnosis, and
regression-aware gated promotion. HarnessFix provides a related precedent
for mapping symptoms to candidate categories~\cite{harnessfix}. Two works reach further in the direction
this thesis takes. \textbf{EMAS}~\cite{emas} refuses to land a revision
until its diagnosis has reproduced across samples --- an admission that
reads evidence rather than score, though not the executed graph. And
\emph{Do Agent Optimizers Compound?}~\cite{compound} asks whether
multi-round gains accumulate at all, which is the repeated-independent-proposal
pattern of \S\ref{sec:archaeology} posed as a benchmark question.

\paragraph{Generation against admission.} Several systems read trace text to propose --- GEPA reads execution traces
and HarnessForge is guided by attribution --- before selecting on score-side
fitness. HarnessBank is the exception: its activation gate requires a
candidate's declared increment to emit a deterministic beacon before paired
significance and gain gates are applied. GHX adds a native execution DAG and
reconciles each candidate's declaration with observed execution.

\paragraph{Evaluation, an orthogonal concern.} A contemporaneous
critique~\cite{rethink-harness-eval} argues that harness-evolution methods
which search and then report on the same public benchmark conflate two
questions the field has not separated --- a concern this thesis's own
evidence argument echoes inside one system rather than across a
leaderboard. None of the harness-level works identified in this survey reports an
independent execution of another group's released loop. Regimes motivates
held-out final evaluation~\cite{regimes}, while Better Harnesses motivates
cross-model checks~\cite{better-harnesses}. This survey did not identify an
earlier external re-execution of the AEGIS release
before Chapter~\ref{ch:baseline}.

% -------------------------------------------------------------------------
\section{Failure attribution: from trace benchmarks to causal representations}
\label{sec:attribution}
% -------------------------------------------------------------------------

\subsection*{Reported limits of trace-based attribution}

Trace-based benchmarks report different endpoints: Who\&When~\cite{whowhen}
reports \textbf{14.2\%} step-level accuracy, while TRAIL~\cite{trail} reports
roughly \textbf{11\%} joint failure classification and localization for its
best model on its 148-trace benchmark. The metrics are not directly
comparable. These results motivate more explicit attribution structure.

\subsection*{Structured attribution}

Later systems introduce explicit causal or provenance-aware representations.
MAP-Graph supplies a provenance- and action-risk-aware gate~\cite{mapgraph}. CHIEF~\cite{chief}
makes the move explicit in its title, and GRADE~\cite{grade} grades edges
as observed, declared or inferred --- a methodological precedent for the
declaration-against-observation reconciliation of \S\ref{sec:unfolded},
and evidence that reconstruction error is a recognized cost in this line.

\subsection*{Attribution linked to repair}

Four works do connect attribution back to a repair action, and the claim
that attribution is ``offline, after the run'' holds for the first wave but
not for them:

\begin{itemize}
  \item \textbf{AgentDebug}~\cite{agentdebug}: detect, attribute, recover,
        re-run.
  \item \textbf{DoVer}~\cite{dover}: intervene at the suspected step,
        re-run, and count single-intervention flips.
  \item \textbf{CausalFlow}~\cite{causalflow}: step-level counterfactual
        responsibility, feeding repair and training pairs.
  \item \textbf{FlowFixer}~\cite{flowfixer}: symbolic root cause to
        pre-execution filter to validated repair, judge-free, reporting
        $71.3\%$ repair success across three platforms.
\end{itemize}

Causal Agent Replay~\cite{car} is deliberately \emph{not} in this set: per
our own reading it stops at diagnosis and was not reproduced in a live
system. It is used here as a methodological precedent
(\S\ref{sec:positioning}).

\subsection*{Comparison across four functions}

Against those four the position separates into four parts:

\begin{enumerate}
  \item \textbf{Native recording.} GHX's unfolded execution graph $U$ is
        written by the runtime, so every edge is observed-grade and none is
        reconstructed --- coverage is bounded by what the runtime observed,
        but no dependency is inferred, the pain GRADE's three-tier grading
        makes explicit --- and the object recorded is the dynamic unfolded execution. The neighbouring
        systems use distinct representations: static orchestration graphs,
        reconstructed dependencies, symbolic traces, or CausalFlow's instrumented
        typed steps.
  \item \textbf{Cross-task aggregation.} Lift, cone diffs and regression
        accounts are computed over a whole bed. Three of the four are
        single-trace diagnoses; FlowFixer aggregates across platforms
        rather than across a bed inside one evolving system.
  \item \textbf{Read--write closure.} The evidence feeds the Evolver,
        passes gates, and lands in the configuration through a
        transactional apply, genotype paired to $U$. None writes its representation back to a persisted harness.
  \item \textbf{Resident intervention.} The edit stays installed and is
        re-judged every round rather than repairing one trajectory once,
        which is testable: it is resident if it survives into round
        $k{+}1$'s configuration.
\end{enumerate}

\paragraph{Comparison with DoVer and FlowFixer.} DoVer intervenes and
re-runs, providing an efficacy-readout pattern, but reconstructs a graph
after one trace, applies a one-off intervention to that trace, and neither
persists the change in a harness nor re-evaluates it in a later round.
FlowFixer provides a closer functional comparison: it is symbolic and judge-free, which
GHX's gates also are, and its $71.3\%$ repair success is a number this
thesis has no equivalent for. What it lacks is residence --- its repair is
validated once and does not survive into a loop that would re-judge it.

A 2026 survey~\cite{survey-attr-evolve} treats attribution and
self-evolution as deeply connected but still-fragmented research threads,
and proposes closing the loop between them.


% -------------------------------------------------------------------------
\section{Execution provenance and AgentOps}
\label{sec:provenance}
% -------------------------------------------------------------------------

A parallel line treats execution records as first-class infrastructure:
OpenTelemetry's \texttt{gen\_ai} semantic conventions, AgentOps
taxonomies, eBPF-based observability.

A 2026 survey of this line~\cite{survey-provenance} organizes execution
provenance along six taxonomy dimensions and describes typed relations that
support verification, attribution, debugging, safety, audit and recovery. It
identifies recovery-oriented evaluation and intervention as open challenges.

Those challenges motivate GHX's design and name this thesis's object: \emph{provenance-grounded} self-evolution, in which admission checks candidate prose against an executed record
instead of relying on prose alone.

% -------------------------------------------------------------------------
\section{Structured candidate spaces}
\label{sec:structured-spaces}
% -------------------------------------------------------------------------

Two systems in this line provide direct methodological precedents for GHX:
MermaidFlow and AgentFlow, taken in turn below, with the relationship to
each stated alongside the design it informed. MEGA's typed evidence atoms
provide a further representation precedent~\cite{mega}.

\paragraph{MermaidFlow: what was taken, what was changed.}
MermaidFlow~\cite{mermaidflow} represents a workflow as a Mermaid graph
over six typed declarative units and evolves it with graph-edit operators
--- crossover, mutation, insert, delete --- under a validator of five
regular-expression checks plus a Mermaid compile. That is enough to lift
the valid-candidate rate above $90\%$ against roughly $50\%$ for AFlow, at
a mean score of $80.75$ against $78.67$. \emph{Taken:} the premise that
the candidate space should be graph edits under a validator rather than
free text, and the operator vocabulary --- GHX's seven typed operations
(\S\ref{sec:candidate-surface}) are that vocabulary specialized to a
processor plane. \emph{Changed:} the object is the harness's composition
layer, not a per-task workflow; the validator is semantic and build-time
rather than syntactic (a write\,$\to$\,reload\,$\to$\,re-graph round trip,
a parity door, a hash-contract test); and identity is treated as a
problem. MermaidFlow's closure lemma states that validated edits compose
``without revalidation'', yet per our mechanism-level read the
implementation re-runs the whole checker on every candidate --- no
incremental revalidation, no graph-hash identity, and no reconciliation of
declaration against observation, there being no execution trace in the
system at all. The ``free of revalidation'' design does not hold at the
composition layer, where an edit to one processor's ordering changes what
every later processor sees.

\paragraph{AgentFlow: what was taken, what was changed.}
AgentFlow's agent dependency graph~\cite{agentflow} recovers, from
heterogeneous agent frameworks, six typed node classes (agent, prompt,
model, capability, memory, policy) and \textbf{three strongly-typed edge
families}: component dependency, guarded control flow, and data flow
including state reads and writes. It is static, LLM-free and an
over-approximation by design --- ``possible dependencies rather than a
single concrete execution trace'' --- recovering a median of 17 nodes and
28 edges on 59 of 60 frameworks. \emph{Taken:} the three edge families,
verbatim, with model-node intent annotation related to AER~\cite{aer} (\S\ref{sec:composition-graph}), and the reaching-definitions
treatment of state reads and writes that gives GHX's data edges their
semantics (\S\ref{sec:unfolded}). \emph{Changed:} the direction of every
one of its design choices. AgentFlow is offline, static and possible;
$U$ is online, concrete and observed, an edge present only if the
dependency it names occurred. AgentFlow's product is a report read by a
person; $U$ is read by an acceptance gate. Its node vocabulary is not
adopted because it describes a different design space. GHX's current
composition exporter materializes hook points, processors, and slots; its
unfolded graph adds processor, tool-call, and nested-agent events observed
in one run. AgentFlow is the static over-approximation; $U$ is GHX's
dynamic, single-trace counterpart.

\paragraph{Relation to evidence-based admission.} This line provides
structural validity constraints for the candidate space: typed units, constraint-preserving
edits, high valid-code rates. The reviewed systems do not connect the typed
candidate to evidence that it \emph{fired}. Typing governs what may be
\emph{proposed}; it never governs what \emph{did} happen.

% -------------------------------------------------------------------------
\section{Graph memory and experience reuse}
\label{sec:graph-memory}
% -------------------------------------------------------------------------

This line is adjacent rather than on the route, and is included to mark the
boundary.


GraphMind~\cite{graphmind} reinforces successful traversal paths in an
action-centric workflow graph, and provides an architectural precedent for
one of GHX's guidance channels. Where this line updates workflow knowledge
and traversal guidance, GHX records execution history for admission.

\paragraph{Boundary.} GraphMind updates workflow knowledge and traversal
guidance; GHX writes persisted harness configuration. Both can influence
future execution, but they update different objects. \S\ref{sec:survivor}
reports that the one channel which produced a positive result in this study
fell on the guidance side of that distinction.

% -------------------------------------------------------------------------
\section{Positioning}
\label{sec:positioning}
% -------------------------------------------------------------------------


\begin{table}[t]
\centering\small
\setlength{\tabcolsep}{5pt}
\renewcommand{\arraystretch}{1.3}
\begin{tabular}{@{}p{2.5cm} p{4.6cm} p{4.9cm}@{}}
\toprule
\textbf{Layer} & \textbf{AEGIS release~\cite{aegis}} & \textbf{GHX (this work)} \\
\midrule
The harness &
  a configuration file read by a runloop; processor order is list order
  plus registration side effects &
  a typed composition graph that \emph{is} the ordering authority, at
  byte-level parity with that runloop \\
A candidate &
  a manifest and a \texttt{config.yaml} composed as prose &
  an atomic group of typed candidate operations, both artifacts
  machine-generated \\
Identity of an edit &
  none downstream: the document and the applied change cannot be
  reconciled &
  graph-layer genotype $\subseteq$ deployment $\subseteq$ phenotype hashes,
  plus generated-config and edit-ledger inspectability for registry operations \\
Evidence about a run &
  a flat log, and a verbatim-overlap heuristic over it, rendered as prose
  dossiers &
  an unfolded execution DAG, observed edges only, with ancestor cones over
  it \\
What admission reads &
  the candidate's prose, at five deterministic gates &
  the prose, \emph{and} the recorded execution, at a sixth \\
The retention decision &
  aggregate score, with a post-round score-drop clamp &
  \emph{unchanged} --- no efficacy readout in either \\
\bottomrule
\end{tabular}
\caption[The released loop against GHX, layer by layer]{The released loop
and GHX, layer by layer. Five layers are re-represented, while the
retention decision remains unchanged, as discussed in
Chapters~\ref{ch:results} and~\ref{ch:discussion}. The release provides
five deterministic gates; GHX adds a sixth check over recorded execution.}
\label{tab:layers}
\end{table}

The contribution claimed here is the implemented integration of four
functions in one resident loop: a native execution graph, cross-task
aggregation over its evidence, write-back to a persisted harness, and
residence across rounds. Table~\ref{tab:layers} sets the released loop
against GHX layer by layer; Table~\ref{tab:position} compares the nearest
systems on those functions. That comparison is a functional positioning aid, not a proof
of priority or an exhaustive novelty census; GHX's typed candidate
operations and graph-layer identity are additional architectural properties
described separately in \S\ref{sec:candidate-surface} and
\S\ref{sec:identity}.

\begin{table}[t]
\centering\small
\setlength{\tabcolsep}{5pt}
\renewcommand{\arraystretch}{1.25}
\begin{tabular}{@{}l cccc@{}}
\toprule
 & \textbf{Native} & \textbf{Cross-task} & \textbf{Write-back to} & \textbf{Resident} \\
\textbf{System} & \textbf{execution graph} & \textbf{aggregation} & \textbf{persisted harness} & \textbf{across rounds} \\
\midrule
HarnessForge~\cite{harnessforge} & --            & \checkmark & \checkmark    & \checkmark \\
AgentDebug~\cite{agentdebug}     & --            & --         & --            & --         \\
DoVer~\cite{dover}               & reconstructed & --         & --            & --         \\
CausalFlow~\cite{causalflow}     & typed trace    & --         & --            & --         \\
FlowFixer~\cite{flowfixer}       & symbolic      & partial    & --            & --         \\
AEGIS release~\cite{aegis}      & --            & prose      & \checkmark    & \checkmark \\
\midrule
\textbf{GHX} (this work)         & \checkmark    & \checkmark & \checkmark    & \checkmark \\
\bottomrule
\end{tabular}
\caption[Nearest systems against four neighbouring functions]{A functional
comparison of the nearest systems against four properties separated by the
counter-examples of \S\ref{sec:attribution}; it is not an exhaustive
novelty proof. A checkmark in the write-back column records persistence,
which is why HarnessForge receives one. HarnessForge writes back an executable Planning--Action--Memory
bundle; GHX additionally supplies typed candidate operations and graph-layer identity, properties discussed in
the text rather than encoded by that column. \emph{Partial} records
FlowFixer's aggregation across platforms rather than across a bed inside
one evolving system. The AEGIS release holds the two right-hand functions;
its execution record is a flat log and its aggregation is prose over
dossiers.}
\label{tab:position}
\end{table}

\subsection*{Differences in related design choices}

Four differences are considered: free-code
candidate spaces~\cite{adas,dgm}, where this thesis argues the candidate
should be typed; natural-language harnesses~\cite{nlah}, where it argues
that a harness whose constraints cannot be machine-checked cannot be
gated, as NLAH's own appendix concedes; knowledge-level
graphs~\cite{mega}, where it argues the graph should be execution-level;
and gate ladders~\cite{alphaevolve,harnessbank}, where it argues that
score significance alone does not establish the mechanism attribution used
to admit a candidate. HarnessBank provides a named comparator through its
activation and significance gates;
the local gain-face analysis shows that score and proxy movements vary
across the recorded events (\S\ref{sec:gain-face}).


